% SGA 3, Expose VII_B, English reconstruction -- Loop 1.
% Scope: Section 5.1.3 in full on authority local pages 85--86;
% stops immediately before Section 5.1.4.
% Combined-reader pages: 598--599. Re-edition printed pages: 587--588.
% Editorial note 148 occurs immediately before this component and is
% preserved in component 78; no editorial note begins inside this component.
% Authority: Polo--Gille Exp7B-13oct24.pdf,
% SHA-256 D303690E0195654BBE7C3DAA5D89EA23E24CD709A3575217234AC19AE8836694.
% Jacob Reinhold's CC BY 4.0 Markdown was consulted as comparison lineage.
% The genuine kernel-pair diagram is an exact PDF-derived Loop 1 crop.

\subsubsection*{5.1.3}
\phantomsection
\label{sga3:VIIB:5.1.3}

将每个预有限 \(k\)-代数 \(\Gamma\) 与满足
\[
  \Gamma_0=\Gamma,
  \qquad
  \Gamma_n=0\quad\text{若 }n>0.
\]
的分次预有限 \(k\)-代数 \((\Gamma_n)_{n\geq0}\) 等同。特别地，若 \((A_n)_{n\geq0}\) 为分次预有限 \(k\)-代数，我们将不加区别地把 \(A_0\) 看作预有限 \(k\)-代数或分次预有限 \(k\)-代数。用
\[
  \rho:(A_n)\longrightarrow A_0
\]
表示 \(\mathbf{Alpg}/k\) 中满足
\[
  \rho_0=\operatorname{id}_{A_0},
  \qquad
  \rho_n=0\quad\text{若 }n>0.
\]
的态射。同样，用
\[
  \tau:A_0\longrightarrow(A_n)
\]
表示 \(\rho\) 的满足
\[
  \tau_0=\operatorname{id}_{A_0},
  \qquad
  \tau_n=0\quad\text{若 }n>0.
\]
的截面。

\((A_n)_{n\in\mathbb N}\) 上的每个余群结构都在 \(A_0\) 上诱导一个余群结构，使 \(\rho\) 和 \(\tau\) 为余群同态。此时，令 \(I_0\) 为 \(A_0\) 的增广理想，并置
\[
  A'_n=A_n/I_0A_n
  \qquad(n\geq0),
\]
于是 \(A'_0=k\)。

则 \((A'_n)_{n\in\mathbb N}\) 为 \(\mathbf{Alpg}/k\) 中的余群。事实上，由于
\[
  A_0=I_0\oplus k\cdot1,
\]
对每个 \(n\)，\(A_n\) 的直和项
\[
  I_0\widehat\otimes_kA_n\simeq I_0A_n
\]
都有定义。由于 \(\tau\) 为 \(\rho\) 的截面，由 \SGARefLink{sga3:VIIB:2.4.A}{2.4.A} 可知余群 \((A_n)_{n\in\mathbb N}\) 为 \(A_0\) 与余群 \((A'_n)_{n\in\mathbb N}\) 的“半直余积”。

更确切地说，\((A'_n)_{n\in\mathbb N}\) 作为 \(\mathbf{Alpg}/k\) 的对象，同构于如下映射对的核：
\begin{center}
  \phantomsection
  \label{sga3:VIIB:5.1.3:graded-cogroup-kernel-pair}
  % SGA3-VII-DIAG-126
  % SGA3-VII-DIAG-126
% Authority: Exp7B-13oct24.pdf, local p. 86.
% Crop: figures/Exp7B_localp086_graded_cogroup_kernel_pair.png
% Locator: tex/components/79_expose_VIIB_section513_graded_cogroup_splitting_en.tex, Loop 1 line 77.
\begin{tikzpicture}[baseline=(lhs.base)]
  \node (lhs) {$(A_n)$};
  \node[right=2.6cm of lhs] (rhs)
    {$(A_n)\widehat\otimes_k A_0$};
  \draw[-{Straight Barb[length=4pt,width=4pt]}]
    ([yshift=2.5pt]lhs.east) --
    node[above] {$\tau_1$}
    ([yshift=2.5pt]rhs.west);
  \draw[-{Straight Barb[length=4pt,width=4pt]}]
    ([yshift=-2.5pt]lhs.east) --
    node[below] {$(\operatorname{id}\widehat\otimes\rho)\Delta$}
    ([yshift=-2.5pt]rhs.west);
\end{tikzpicture}

\end{center}
这里
\[
  \Delta:(A_n)\longrightarrow(A_n)\widehat\otimes(A_n)
\]
为 \((A_n)\) 的余乘法，且
\[
  \tau_1(x)=x\widehat\otimes1.
\]
将 \(A'_n\) 与其在 \(A_n\) 中的像等同，则映射
\[
  A'_n\widehat\otimes_kA_0\longrightarrow A_n,
  \qquad
  a'_n\widehat\otimes a_0\longmapsto a'_na_0,
\]
为 \(\mathbf{Alpg}/k\) 中的同构。

这不是余群同构。实际上，有
\[
  \Delta(A')\subset A\widehat\otimes A',
\]
及
\[
  \bigl(\gamma'_\rho\widehat\otimes\operatorname{id}\bigr)
  \circ\Delta\big|_{A'}
  =
  \Delta',
\]
其中 \(\Delta'\) 为 \(A'\) 的余乘法，\(\gamma'_\rho:A\to A'\) 为投影\textup{（参见 \SGARefLink{sga3:VIIB:2.4.A}{2.4.A}）}。
